\section{Розділ~\ref{sec:chemoutput}. Хімічний вихід}

Два проводи до серця та їхню різну швидкість, від'ємний зворотний зв'язок у гормональних каскадах, код частотою порцій, добовий генератор і його підстроювання світлом виміряно прямо; межа $K<8$ --- теорема теорії автоматичного керування, виведена в розділі; пульсації, породжувані самим зворотним зв'язком каскаду, --- гіпотеза із сильним дослідом на її користь.

{\footnotesize
\setlength{\tabcolsep}{3pt}\renewcommand{\arraystretch}{1.15}
\begin{longtable}{>{\raggedright}p{0.5cm}>{\raggedright}p{4.0cm}>{\raggedright}p{5.6cm}>{\raggedright}p{2.3cm}>{\raggedright\arraybackslash}p{1.7cm}}
\toprule
№ & Твердження & Дослід і що виміряно & Джерело & Статус \\
\midrule\endfirsthead
\toprule
№ & Твердження & Дослід і що виміряно & Джерело & Статус \\
\midrule\endhead
1 & Ритмоводій серця сам б'ється близько $100$ разів на хвилину; у спокої гальмо натиснуте, і частота зазвичай порядку $60$--$70$ & Люди, повне вимкнення обох проводів до серця: власна частота вища за частоту спокою і спадає з віком, у дорослих порядку $100$ ударів на хвилину. Отже, у спокої переважає гальмо. Частота спокою $60$--$70$ --- типове значення, окремо не цитоване. & \cite{JoseCollison1970ihr} & виміряно \\
2 & Газ \nt{NORA} і гальмо \nt{ACET} --- фільтри нижніх частот, гальмо швидше~\eqref{eq:gasbrake} & Собака, частотно-модульована стимуляція блукаючого або симпатичного нерва серця й передатна функція до частоти передсердь: обидва шляхи --- фільтри нижніх частот, у симпатичного частота зрізу значно нижча й додано чисту затримку близько $1.7\,\text{с}$. Характеристики залежать від середнього тонусу, тож лінійна формула~\eqref{eq:gasbrake} --- наближення. & \cite{Berger1989} & виміряно \\
3 & Гальмо швидке, бо калієвий канал \nt{ACET} відкриває без другого посередника, а \nt{NORA} --- через ланцюжок реакцій & Мембранні клаптики клітин передсердь: калієвий канал, що відкривається ацетилхоліном, відкривають $\beta\gamma$-субодиниці G-білка, пов'язаного з рецептором, без другого посередника всередині клітини. Шлях \nt{NORA} через цАМФ тут не цитовано. & \cite{Logothetis1987gbg} & виміряно \\
4 & Кров обходить тіло за час порядку хвилини; \nt{ADRE} з кровотоку доходить до всіх органів із рецептором & Загальна оцінка порядку величини; у розділі так і сформульовано, джерело не цитовано. & --- & оцінка \\
5 & Гормональний каскад охоплено від'ємним зворотним зв'язком: кінцевий гормон гальмує перші ступені & Огляд дослідів: гормони кори надниркових залоз пригнічують викид АКТГ гіпофізом і рилізинг-гормону гіпоталамусом у трьох часових діапазонах --- швидкому, проміжному й повільному. & \cite{KellerWood1984fb} & виміряно \\
6 & Три однакові ступені стійкі за $K<8$~\eqref{eq:cascadestable}, на межі період $2\pi\tau/\sqrt3\approx3.6\tau$ & Виведено в розділі: на частоті $\sqrt3/\tau$ три ступені дають зсув фази $180^\circ$ і ослаблення у $8$ разів. & виведення в розділі & теорія \\
7 & Похибка рівня $1/(1+K)$, тому точніше приблизно за $10\,\%$ такий регулятор не тримає (твердження~\ref{prop:cascade}) & Наслідок рядка~6 для пропорційного зворотного зв'язку. Справжня вісь має зворотний зв'язок на кількох ступенях і в кількох часових діапазонах (рядок~5), тож межа стосується моделі~\eqref{eq:cascade}, а не виміряна. & виведення в розділі & теорія \\
8 & За $K=4$ і $7$ рівень усталюється, за $K=12$ --- рівні пульсації з періодом $3.7$--$3.9\,\tau$ (рис.~\ref{fig:cascade}) & Розрахунок за \texttt{models/\allowbreak chem\_\allowbreak models.py}: за $K=12$ послідовні періоди $3.9$, $3.7$, $3.8$, $3.7\,\tau$; за $K=4$ розмах наприкінці рахунку $0.003$. & \texttt{models/\allowbreak chem\_\allowbreak models.py} & модель \\
9 & Гормон стресу виходить пульсами приблизно раз на годину; пульси породжує сам зворотний зв'язок, без окремого генератора & Щури, постійне вливання рилізинг-гормону: АКТГ і кортикостерон коливаються з фізіологічною, майже годинною частотою --- ритмічний вхід від гіпоталамуса для пульсів не потрібен. Модель гіпофіз--наднирник із запізнілим зворотним зв'язком дає такі коливання. & \cite{Walker2012glc, Walker2010} & гіпотеза \\
10 & Отримувач читає частоту порцій, а не рівень & Мавпи без власного викиду рилізинг-гормону гонадотропінів: постійне вливання не відновлює викиду гормонів гіпофіза, порції раз на годину відновлюють; перехід на постійне вливання знову глушить відповідь. Утома рецептора як фільтр верхніх частот --- тлумачення. & \cite{Belchetz1978} & виміряно \\
11 & Різні частоти порцій по-різному діють на різні клітини однієї залози & Ті самі мавпи: почастішання порцій до $2$--$5$ на годину знижує обидва гормони гіпофіза; порідшання до однієї за $3$~години знижує один (ЛГ) і підвищує другий (ФСГ), тож їхнє відношення задається частотою. & \cite{Wildt1981gnrh} & виміряно \\
12 & Добовий ритм породжує внутрішньоклітинний від'ємний зворотний зв'язок із затримкою в кілька годин & Огляд: білки годинникових генів із затримкою пригнічують транскрипцію власних генів; петля з транскрипції й трансляції дає період близько доби. & \cite{TakahashiJS2017} & виміряно \\
13 & Головний генератор --- група з десятків тисяч нейронів, що синхронізує решту тіла & Огляд: супрахіазматичне ядро --- головний добовий генератор; окремі його нейрони в культурі --- самостійні генератори, мережа їх синхронізує; у миші близько $10^4$ нейронів із кожного боку. & \cite{Welsh2010scn} & виміряно \\
14 & Власний період у людини $24.18$~години & Люди в тьмяному світлі за розкладом, відірваним від доби: період ритмів мелатоніну, температури й кортизолу в середньому $24.18$~години в молодих і літніх, із вузьким розкидом. & \cite{Czeisler1999} & виміряно \\
15 & Світло підводить генератор як фазове автопідстроювання~\eqref{eq:pll}; потрібен зсув близько $11$~хв на добу & Люди, одиночне яскраве світло $6.7$~години в різний час доби: крива фазової відповіді першого типу з розмахом $5$~годин --- затримка до мінімуму температури тіла, випередження після. Рівняння~\eqref{eq:pll} --- стандартна модель; $11$~хв $=0.18$~години --- арифметика. & \cite{Khalsa2003prc} & виміряно \\
16 & Без світла підстроювання немає, і ритм переходить на власний період & Цілком незрячі люди без сприйняття світла, що живуть у звичайному суспільстві: приблизно в половини ритми мелатоніну й кортизолу йдуть зі своїм періодом, що не збігається з добою. & \cite{Sack1992blind} & виміряно \\
\bottomrule
\end{longtable}}
